\scanpage{187}
\mspage{500}
% position: top continuation
tant bien que mal. Cette réalité ``sociale'' est le produit d'un conditionnement
immémorial, lequel très tôt prend racine dans le ``moi'' en formation et le struc-
ture. Pourtant, au delà de cette réalité-là, il est une réalité plus profonde,
venant de beaucoup plus loin encore, laquelle est déterminante dans la pulsion
amoureuse elle-même. C'est la réalité d'une \mbox{c o m p l é m e n t a r i t é}
profonde, essentielle des sexes, où il n'y a nulle place pour un quelconque
``antagonisme''. C'est la réalité aussi qui se manifeste clairement dans toutes les
espèces vivantes, à la seule exception de la nôtre, où elle se trouve occultée
dans une large mesure par l'antagonisme culturel, donc par un état de \mbox{d i v i-}
\mbox{s i o n} propre à l'homme et à la société humaine.

% position: long middle paragraph
Les clichés courants romantiques, style ``Nous Deux'', qui dominent une
grande partie de la littérature et des médias, montent d'ailleurs en épingle
une ``complémentarité'' de pacotille, tout en jetant un voile pudique sur le trou-
blant aspect antagoniste homme-femme, ou (au mieux) en le traitant comme une
sorte d'accident un peu piquant, bienvenu pour mettre quelque piment dans un
repas un peu trop fadasse ou sirupeux sinon. Dès qu'on dépasse ce gen\inserted{re de}
clichés rassurants, on se voit aussitôt confronté à la réalité de cet antagonis-
me homme-femme - réalité apparemment universelle, et au surplus d'une tenacité à
toute épreuve, une tenacité de chiendent! Mais \quad partir de cette réalité omni-
présente et irrécusable, pour instituer une sorte d'antagonisme cosmique du
yin et du yang, du ``féminin'' et du ``masculin'', c'est projeter sur l'Univers
entier l'état de déchirement, de division profonde de la société humaine et de
la personne, une maladie donc propre à notre espèce. C'est aussi perpétuer sa
propre ignorance d'une \mbox{a u t r e} réalité en soi-même (rejoignant cette
réalité cosmique de l'harmonie des complémentaires), d'une réalité toute aussi
tenace (ou, pour mieux dire, indestructible), mais plus cachée. Cette réalité
va à l'encontre des conditionnements instituant tacitement un antagonisme de
fait aussi bien entre la femme et l'homme, l'épouse et l'époux, qu'entre cela
en nous-même qui est ``femme'' et ce qui est ``homme''.

% position: lower paragraph
A vrai dire, cette vision \mbox{d u a l i s t e} ou \mbox{g u e r r i è r e} de
l'Univers, où un aspect des choses se trouverait en guerre constante avec un
aspect ``symétrique'' tout aussi essentiel - cette vision n'est nullement le
fruit d'une \mbox{r é f l e x i o n}, qui ``partirait'' (comme j'écrivais à l'instant)
de la réalité du conflit dans le couple humain et dans la société humaine, pour
la ``déduire'' ensuite (ou ``l'instituer'', comme j'écrivais plus justement) dans
le Cosmos tout entier. Elle n'est ni plus, ni moins que l'expression fidèle,
automatique autant dire, du conditionnement culturel, et va dans le sens d'une
\margin{-}
\margin{186}
